\documentclass[conference]{IEEEtran}

\usepackage[utf8]{inputenc}
\usepackage[T1]{fontenc}
\usepackage{cite}
\usepackage{amsmath,amssymb,amsfonts}
\usepackage{booktabs}
\usepackage{graphicx}
\usepackage{hyperref}
\usepackage{xcolor}

\hypersetup{colorlinks=true, linkcolor=blue, citecolor=blue, urlcolor=blue}

\begin{document}

\title{Fine-Tuning Small LLMs for Turkish Legal Question Answering:\\
A QLoRA-Based Approach}

\author{
  \IEEEauthorblockN{Group 8 --- CSE4078 Spring 2026}
  \IEEEauthorblockA{
    Ayşegül Bihter Banuşoğlu (150119777) \quad Kaan Camcı (150119063) \\
    Yakup Mert Aslan (150119765) \quad Muhammed Taha Serdar (150122604) \quad Muhammet Akyüz (150120028)
  }
}

\maketitle

% ---------------------------------------------------------------
\begin{abstract}
This paper presents an evaluation and fine-tuning study of small open-source large language
models (LLMs) for Turkish legal question answering. We evaluate two models in the 1B--4B
parameter range---Qwen2.5-1.5B-Instruct and Llama-3.2-3B-Instruct---as baselines on the
\textit{Renicames/turkish-law-chatbot} dataset. We then fine-tune Llama-3.2-3B-Instruct
using Quantized Low-Rank Adaptation (QLoRA) on the full 13,354-example training split
for five epochs. Evaluation on the 1,500-example test split using seven automatic metrics
(ROUGE-1/2/L, BLEU, chrF++, TF-IDF similarity, and BERTScore F1) shows substantial
improvements across all metrics after fine-tuning. Our fine-tuned model achieves a ROUGE-L
of 0.5874 compared to 0.2506 for the Llama baseline and 0.2126 for Qwen, demonstrating that
lightweight parameter-efficient fine-tuning can significantly close the gap between
general-purpose and domain-specific legal QA in Turkish.
\end{abstract}

% ---------------------------------------------------------------
\section{Introduction}

Turkish is a morphologically rich, agglutinative language that poses unique challenges for
natural language processing systems. Despite recent advances in multilingual LLMs, Turkish
legal question answering remains an underexplored task. Legal texts in Turkish are
particularly demanding due to formal register, complex sentence structures, and
domain-specific vocabulary.

Recent work has shown that parameter-efficient fine-tuning methods such as
LoRA~\cite{hu2022lora} and QLoRA~\cite{dettmers2023qlora} can adapt general-purpose LLMs to
specialized domains without the computational cost of full fine-tuning. These methods are
especially attractive for small models in the 1B--4B parameter range that can be trained
on consumer GPUs.

In this work, we conduct a systematic study of two small open-source LLMs on Turkish legal
QA. Our contributions are:
\begin{itemize}
  \item A comparative baseline evaluation of Qwen2.5-1.5B-Instruct and
    Llama-3.2-3B-Instruct on a Turkish legal QA benchmark.
  \item A QLoRA fine-tuning pipeline for Llama-3.2-3B-Instruct on the
    \textit{turkish-law-chatbot} corpus.
  \item A comprehensive evaluation using seven automatic metrics on a 1,500-example test
    split, providing a reproducible benchmark for future work.
  \item An error analysis comparing model outputs before and after fine-tuning.
\end{itemize}

\subsection{Related Work}

Instruction-tuned LLMs such as LLaMA~\cite{touvron2023llama} and Qwen~\cite{bai2023qwen}
have demonstrated strong multilingual capability. However, their Turkish legal domain
performance without explicit fine-tuning remains limited.
LoRA~\cite{hu2022lora} introduced low-rank decomposition of weight matrices for
parameter-efficient fine-tuning. QLoRA~\cite{dettmers2023qlora} extended this approach by
combining 4-bit quantization with LoRA, enabling fine-tuning of large models on single
consumer GPUs. We follow the QLoRA approach in this work.
Multilingual evaluation metrics such as BERTScore~\cite{zhang2020bertscore} with
XLM-RoBERTa~\cite{conneau2020xlmroberta} and character-level metrics like
chrF++~\cite{popovic2017chrf} are more suitable for agglutinative languages like Turkish
than standard BLEU, and we include both in our evaluation.

% ---------------------------------------------------------------
\section{Experiment Setup}

\subsection{Dataset}

We use the \textit{Renicames/turkish-law-chatbot} dataset~\cite{turkishlawchatbot}, which
contains 14,900 Turkish legal question-answer pairs. The dataset is split into a training
set of 13,354 examples and a test set of 1,500 examples under the Apache-2.0 license.
Each example consists of a question (\textit{Soru}) and a reference answer (\textit{Cevap})
in Turkish. The training split was used exclusively for fine-tuning; the test split was
held out and used only for evaluation.

We applied no additional preprocessing beyond the chat template formatting described below.
We verified that no examples from the test split appear in the training data, ensuring a
clean train/test separation.

\subsection{Baseline Models}

We selected two instruction-tuned models in the 1B--4B parameter range:

\begin{itemize}
  \item \textbf{Qwen2.5-1.5B-Instruct}~\cite{bai2023qwen}: A 1.5B parameter model from
    the Qwen2.5 series, optimized for instruction following and multilingual tasks.
  \item \textbf{Llama-3.2-3B-Instruct}~\cite{touvron2023llama}: A 3B parameter model
    from Meta's Llama 3.2 series with strong multilingual performance.
\end{itemize}

Both models were loaded in BFloat16 precision with \texttt{device\_map="auto"} for GPU
inference. Prompts were formatted using each model's native chat template with a shared
system prompt instructing the model to answer Turkish legal questions concisely and
accurately.

\subsection{Fine-Tuning}

We selected Llama-3.2-3B-Instruct for fine-tuning based on its stronger baseline
performance. Fine-tuning was performed using QLoRA~\cite{dettmers2023qlora} with the
following configuration:

\begin{itemize}
  \item \textbf{Quantization}: 4-bit NF4 with double quantization
    (\texttt{BitsAndBytesConfig})
  \item \textbf{LoRA rank}: $r = 16$, $\alpha = 32$, dropout $= 0.05$
  \item \textbf{Target modules}: \texttt{q\_proj}, \texttt{k\_proj}, \texttt{v\_proj},
    \texttt{o\_proj}, \texttt{gate\_proj}, \texttt{up\_proj}, \texttt{down\_proj}
  \item \textbf{Optimizer}: AdamW with cosine learning rate schedule
  \item \textbf{Learning rate}: $2 \times 10^{-4}$, warmup ratio $= 0.03$
  \item \textbf{Batch size}: 4, gradient accumulation steps: 4 (effective batch size 16)
  \item \textbf{Epochs}: 5
  \item \textbf{Max sequence length}: 512 tokens
  \item \textbf{Precision}: BFloat16
\end{itemize}

Training examples were formatted using the Llama chat template with a Turkish system prompt:
\textit{``Aşağıdaki Türkçe hukuk sorusunu kısa, açık ve doğru şekilde cevapla.''} Training
was performed using the TRL library~\cite{vonwerra2022trl} (\texttt{SFTTrainer}) on an
NVIDIA RTX 3090 (24 GB VRAM). The total training time was approximately 5 hours.

\subsection{Evaluation}

All models were evaluated on the full 1,500-example test split using seven automatic metrics:

\begin{itemize}
  \item \textbf{ROUGE-1, ROUGE-2, ROUGE-L}~\cite{lin2004rouge}: Unigram overlap, bigram
    overlap, and longest common subsequence between prediction and reference.
  \item \textbf{BLEU}~\cite{papineni2002bleu}: N-gram precision with effective-order
    brevity penalty (via SacreBLEU).
  \item \textbf{chrF++}~\cite{popovic2017chrf}: Character n-gram F-score with word order
    (\texttt{word\_order=2}), well-suited for agglutinative languages such as Turkish.
  \item \textbf{TF-IDF Cosine Similarity}: Cosine similarity between TF-IDF vectors of
    prediction and reference, measuring lexical overlap at the corpus level.
  \item \textbf{BERTScore F1}~\cite{zhang2020bertscore}: Contextual embedding similarity
    using XLM-RoBERTa-base with \texttt{lang="tr"}, capturing semantic similarity beyond
    exact token matches.
\end{itemize}

All text was lowercased and whitespace-normalized before metric computation. Generation was
performed greedily (\texttt{do\_sample=False}) with a maximum of 100 new tokens.

% ---------------------------------------------------------------
\section{Experiment Results}

\subsection{Quantitative Results}

Table~\ref{tab:results} presents the evaluation results for all three models on the
1,500-example test split.

\begin{table*}[t]
\centering
\caption{Evaluation results on the Turkish Legal QA test set (1,500 examples).
All scores are averages over the test split. Higher is better for all metrics.}
\label{tab:results}
\begin{tabular}{lccccccc}
\toprule
\textbf{Model} & \textbf{ROUGE-1} & \textbf{ROUGE-2} & \textbf{ROUGE-L} &
\textbf{BLEU} & \textbf{chrF++} & \textbf{TF-IDF} & \textbf{BERTScore F1} \\
\midrule
Qwen2.5-1.5B-Instruct (baseline)     & 0.2609 & 0.1271 & 0.2126 & 0.0688 & 0.2810 & 0.1648 & 0.8667 \\
Llama-3.2-3B-Instruct (baseline)     & 0.2896 & 0.1597 & 0.2506 & 0.1002 & 0.2802 & 0.1979 & 0.8740 \\
Llama-3.2-3B-Instruct (fine-tuned)   & \textbf{0.6186} & \textbf{0.5163} & \textbf{0.5874} & \textbf{0.4720} & \textbf{0.6038} & \textbf{0.5419} & \textbf{0.9369} \\
\bottomrule
\end{tabular}
\end{table*}

The fine-tuned model substantially outperforms both baselines across all metrics. Compared
to the Llama-3.2-3B baseline, fine-tuning yields a ROUGE-L improvement from 0.2506 to
0.5874 (+134.4\%), a BLEU improvement from 0.1002 to 0.4720 (+371.1\%), and a BERTScore
F1 improvement from 0.8740 to 0.9369 (+7.2 points).

The two baseline models perform comparably, with Llama-3.2-3B-Instruct holding a small
advantage on most metrics despite having twice the parameters of Qwen2.5-1.5B-Instruct.
This is consistent with the expectation that the larger model has better instruction
following and Turkish language capability.

\subsection{Error Analysis}

To better understand model behavior, we examined representative outputs from each model.

\textbf{Baseline models} tend to produce plausible but generic answers that do not match
the specific legal facts in the reference answer. For example, when asked about the
conditions for filing a lawsuit, baseline models often give general procedural descriptions
while the reference answer cites specific articles of Turkish law. Both baseline models
sometimes respond with structured bullet points or introductory phrases that are not present
in the reference, lowering n-gram overlap metrics even when the semantic content is broadly
correct.

\textbf{The fine-tuned model} produces answers that are stylistically and semantically much
closer to the reference answers. It correctly reproduces specific legal terminology, article
references, and answer lengths consistent with the dataset distribution. The main failure
modes observed after fine-tuning are: (1) occasional over-repetition of key phrases when
the question is ambiguous, and (2) answers that are slightly longer than the reference,
which mildly penalizes precision-based metrics.

The high BERTScore F1 of the baseline models (0.8667--0.8740) relative to their low BLEU
scores (0.0688--0.1002) confirms that the baselines understand the semantic topic but
cannot reproduce the exact vocabulary and phrasing expected by the dataset. Fine-tuning
closes this gap dramatically.

% ---------------------------------------------------------------
\section{Conclusion}

We presented a study of QLoRA-based fine-tuning for Turkish legal question answering.
Starting from two competitive small LLMs---Qwen2.5-1.5B-Instruct and
Llama-3.2-3B-Instruct---we showed that five epochs of QLoRA fine-tuning on 13,354 examples
from the \textit{turkish-law-chatbot} corpus yields substantial gains across seven automatic
metrics. The fine-tuned Llama-3.2-3B model achieves a ROUGE-L of 0.5874 and a BLEU of
0.4720 on the 1,500-example test split, compared to 0.2506 and 0.1002 for the untuned
baseline.

These results confirm that lightweight parameter-efficient fine-tuning is highly effective
for domain adaptation of small LLMs to Turkish legal QA. Future work could explore longer
training, data augmentation, or retrieval-augmented generation to further improve factual
accuracy on complex legal questions.

% ---------------------------------------------------------------
\begin{thebibliography}{00}

\bibitem{hu2022lora}
E. J. Hu, Y. Shen, P. Wallis, Z. Allen-Zhu, Y. Li, S. Wang, L. Wang, and W. Chen,
``LoRA: Low-Rank Adaptation of Large Language Models,''
\textit{Proc. ICLR}, 2022.

\bibitem{dettmers2023qlora}
T. Dettmers, A. Pagnoni, A. Holtzman, and L. Zettlemoyer,
``QLoRA: Efficient Finetuning of Quantized LLMs,''
\textit{Proc. NeurIPS}, 2023.

\bibitem{touvron2023llama}
H. Touvron et al.,
``Llama 2: Open Foundation and Fine-Tuned Chat Models,''
\textit{arXiv preprint arXiv:2307.09288}, 2023.

\bibitem{bai2023qwen}
J. Bai et al.,
``Qwen Technical Report,''
\textit{arXiv preprint arXiv:2309.16609}, 2023.

\bibitem{turkishlawchatbot}
Renicames,
``turkish-law-chatbot,''
Hugging Face Datasets, 2023.
[Online]. Available: \url{https://huggingface.co/datasets/Renicames/turkish-law-chatbot}

\bibitem{lin2004rouge}
C.-Y. Lin,
``ROUGE: A Package for Automatic Evaluation of Summaries,''
\textit{Proc. ACL Workshop}, 2004.

\bibitem{papineni2002bleu}
K. Papineni, S. Roukos, T. Ward, and W. Zhu,
``BLEU: a Method for Automatic Evaluation of Machine Translation,''
\textit{Proc. ACL}, 2002.

\bibitem{popovic2017chrf}
M. Popovic,
``chrF++: words helping character n-grams,''
\textit{Proc. WMT}, 2017.

\bibitem{zhang2020bertscore}
T. Zhang, V. Kishore, F. Wu, K. Q. Weinberger, and Y. Artzi,
``BERTScore: Evaluating Text Generation with BERT,''
\textit{Proc. ICLR}, 2020.

\bibitem{conneau2020xlmroberta}
A. Conneau et al.,
``Unsupervised Cross-lingual Representation Learning at Scale,''
\textit{Proc. ACL}, 2020.

\bibitem{vonwerra2022trl}
L. von Werra et al.,
``TRL: Transformer Reinforcement Learning,''
GitHub, 2022.
[Online]. Available: \url{https://github.com/huggingface/trl}

\end{thebibliography}

\end{document}
